\chapter{Introdução}
\label{cap:introducao}

% Alinhado a: Chris (Eng.ish) Research Document Structure --- I. Introduction
% Design Science Research (Hevner et al.): Relevance / Design / Rigor cycles.

\section{Contexto e motivação}
\label{sec:contexto}

O CET-ADS articula projetos do Air Domain Study (ADS) no ecossistema CONCEPTIO.
O ADS-VD (\emph{Virtual Demonstrator}) espera, entre outros resultados, um ambiente
que permita a medição de métricas desejadas e a avaliação cruzada de projetos.

Na prática, demos acadêmicos multiagente ainda sofrem de um gap recorrente:
requisitos pouco formalizados antes da run; telemetria rica durante a demo sem
laudo auditável depois; e ausência de um ponto de partida MBSE (esqueleto
SysML~v2) ligado ao que foi validado~\cite{sis08methods}.

A exploração de simulação SysML em ferramentas comerciais mostrou limite para
refletir telemetria live, motivando uma prova de conceito \emph{standalone}
(ReqValLive) que preserve a disciplina de Success Criteria.

\section{Problema de pesquisa}
\label{sec:problema}

Como viabilizar, em escala de laboratório acadêmico (CONCEPTIO), um ciclo
operacional em que requisitos com Success Criteria sejam aprovados antes da run,
a telemetria real de drones via MQTT seja avaliada durante missões de teste, e
ao final se obtenha um laudo PASS/FAIL com momentos críticos mais um esqueleto
SysML~v2 reutilizável --- com extensibilidade natural a eventos de integração
CET-ADS (SARITA) sob protocolo MQTT padronizado?

\subsection{Perguntas derivadas}

\begin{enumerate}
  \item Que regras tornam um Success Criterion aceitável para medição live
        (gate ACCEPT/REJECT)?
  \item Como garantir evidência temporal íntegra (ex.: \texttt{all\_timesteps}),
        e não apenas o último valor amostrado?
  \item O ciclo \emph{requisito $\rightarrow$ SC aprovado $\rightarrow$ MQTT $\rightarrow$
        laudo + esqueleto SysML} é realizável ponta a ponta no CONCEPTIO?
\end{enumerate}

\section{Hipótese}
\label{sec:hipotese}

Se Success Criteria forem (i)~tornados verificáveis e aprováveis \emph{antes}
da run e (ii)~avaliados contra telemetria real de drones com regras temporais
explícitas e protocolo MQTT conhecido, então é possível produzir, ao final de
missões de teste no laboratório, um laudo e um esqueleto SysML que fechem o ciclo
\emph{métrica de sucesso $\rightarrow$ evidência $\rightarrow$ ok/nok} de forma mais
auditável do que demos ad~hoc sem critério pré-declarado.

\section{Objetivos}
\label{sec:objetivos}

\subsection{Objetivo geral}

Propor e demonstrar um método e uma prova de conceito (ReqValLive) para
validação live de Success Criteria em missões multiagente de laboratório,
alinhado a práticas de V\&V e ao papel de medição de métricas no ADS-VD / CET-ADS.

\subsection{Objetivos específicos}

\begin{enumerate}
  \item Formalizar um gate de Success Criteria (ACCEPT/REJECT) antes da medição,
        inspirado em SIS-08 Methods / MSFC-HDBK-3173~\cite{msfc3173}.
  \item Implementar avaliação live de critérios numéricos
        (\texttt{threshold}/\texttt{range}/\texttt{statistical}/\texttt{rate}/\texttt{dwell}) contra telemetria
        MQTT multi-entidade, com evidência temporal íntegra.
  \item Gerar artefatos pós-missão: laudo (PASS/FAIL, momentos críticos) e
        esqueleto SysML~v2 textual~\cite{sysmlv2}.
  \item Demonstrar o ciclo ponta a ponta no CONCEPTIO (obrigatório); SARITA como
        extensão desejável sob payload MQTT uniforme.
\end{enumerate}

\section{Metodologia de pesquisa --- Design Science Research}
\label{sec:dsr}

Este trabalho adopta \emph{Design Science Research} (DSR)~\cite{hevner2004},
com três ciclos:

\begin{description}
  \item[Relevance] Necessidade de medição auditável de Success Criteria em demos
        multiagente no CONCEPTIO / CET-ADS (ADS-VD).
  \item[Design] Artefacto ReqValLive + método (gate, evidência live, laudo,
        esqueleto SysML~v2 / fecho no SysON).
  \item[Rigor] Fundamentação em V\&V / Success Criteria, MBSE SysML~v2, telemetria
        MQTT de laboratório; avaliação por demonstração e casos controlados.
\end{description}

A avaliação do artefacto combina testes de regressão, ensaios com telemetria
real (e fixtures) e demonstração do loop fechado modelo~$\leftrightarrow$~medição.

\section{Organização do documento}
\label{sec:estrutura}

\begin{itemize}
  \item Capítulo~\ref{cap:literatura} --- revisão da literatura (ciclo de rigor)
        e, opcionalmente, estado da arte do domínio (ciclo de relevância).
  \item Capítulo~\ref{cap:metodologia} --- materiais, desenho e uso do artefacto
        (fusão método + PoC).
  \item Capítulo~\ref{cap:resultados} --- resultados e discussão (ciclo de design).
  \item Capítulo~\ref{cap:conclusao} --- conclusões e trabalhos futuros.
\end{itemize}
